\documentclass[11pt,a4paper]{article}

\usepackage[T1]{fontenc}
\usepackage[utf8]{inputenc}
\usepackage{lmodern}
\usepackage[a4paper,top=2.4cm,bottom=2.2cm,left=2.2cm,right=2.2cm]{geometry}
\usepackage{amsmath}
\usepackage{amssymb}
\usepackage{graphicx}
\graphicspath{{./}}
\usepackage{array}
\usepackage{makecell}
\renewcommand{\arraystretch}{1.25}
\usepackage{enumitem}
\usepackage{fancyhdr}
\usepackage{xcolor}

\newcommand{\ohms}{\ensuremath{\Omega}}
\newcommand{\degC}{\ensuremath{^{\circ}}C}
\newcommand{\dg}{\ensuremath{^{\circ}}}
\newcommand{\pow}[1]{\ensuremath{^{#1}}}

\newcommand{\topiclabel}[1]{{\hfill\normalfont\small\itshape\textcolor{black!55}{#1}}}

% Per-question head: \examq{paper-id}{number}{topics}{marks}
% Renders: paper-id - Q<number> - topics - marks  (one line, same font).
% Same command and same rendered line as the answer paper.
\newcommand{\swmarks}{}
\newcommand{\examq}[4]{%
  \gdef\swmarks{#4}%
  \par\addvspace{16pt}\noindent
  {\large\bfseries #1 - Q#2 - #3 - #4}\par\vspace{5pt}%
}
% \total now takes NO argument: it prints the marks stored by \examq.
\newcommand{\total}{\par\vspace{2mm}\noindent\hspace*{\fill}[Total:\ \swmarks]\par\vspace{2mm}}

\newlength{\anslineskip}
\setlength{\anslineskip}{8mm}

\newcounter{ansln}
\newcommand{\Alines}[2]{%
  \par\vspace{1mm}%
  \setcounter{ansln}{1}%
  \loop\ifnum\value{ansln}<#1\relax
    \noindent\mbox{}\dotfill\par\vspace{\anslineskip}%
    \stepcounter{ansln}%
  \repeat
  \noindent\mbox{}\dotfill\hspace{0.6em}\mbox{#2}\par
  \vspace{1mm}%
}

\newcommand{\plainline}{\par\noindent\mbox{}\dotfill\par\vspace{\anslineskip}}

\newcommand{\labelline}[1]{\par\vspace{1mm}\noindent #1\hspace{0.5em}\dotfill\par\vspace{\anslineskip}}

\newcommand{\ansval}[3]{%
  \par\vspace{5mm}%
  \noindent\hspace*{3.5cm}#1\hspace{0.5em}\dotfill\hspace{0.5em}#2\hspace{0.5em}\makebox[1.1cm][r]{#3}\par
  \vspace{3mm}%
}

\renewcommand{\markright}[1]{\par\nopagebreak\noindent\hspace*{\fill}\mbox{#1}\par\vspace{1mm}}

\newcommand{\qfig}[3][0.5\linewidth]{%
  \par\vspace{2mm}
  \begin{center}%
    \IfFileExists{#2}%
      {\includegraphics[width=#1]{#2}}%
      {\fbox{\parbox[c][26mm][c]{#1}{\centering\ttfamily\footnotesize #2}}}\\[4pt]
    \textbf{#3}%
  \end{center}%
  \vspace{2mm}\par
}

\newlist{parts}{enumerate}{1}
\setlist[parts]{label=\textbf{(\alph*)},leftmargin=2.1em,labelsep=0.6em,itemsep=10pt,topsep=6pt,parsep=4pt}

\newlist{subparts}{enumerate}{1}
\setlist[subparts]{label=\textbf{(\roman*)},leftmargin=2.3em,labelsep=0.6em,itemsep=10pt,topsep=6pt,parsep=4pt}

\pagestyle{fancy}
\fancyhf{}
\fancyhead[L]{\footnotesize 5054/11/M/J/25}
\fancyhead[C]{\footnotesize Cambridge O Level Physics -- Paper 1 Multiple Choice}
\fancyhead[R]{\footnotesize May/June 2025}
\fancyfoot[C]{\footnotesize\thepage}
\renewcommand{\headrulewidth}{0.4pt}
\renewcommand{\footrulewidth}{0pt}

\setlength{\parindent}{0pt}

% option list for stacked A--D choices
\newlist{choices}{enumerate}{1}
\setlist[choices]{label=\textbf{\Alph*},leftmargin=2.1em,labelsep=0.9em,itemsep=2pt,topsep=5pt,parsep=0pt}

\begin{document}

%------------------------------ 1 ------------------------------
\examq{5054/11/M/J/25}{1}{Physical Quantities and Measurements}{1}
Which quantity has both magnitude and direction, and is it a scalar quantity or a vector quantity?

\begin{center}
\begin{tabular}{|c|c|c|}
\hline
 & \makecell{quantity\\with magnitude\\and direction} & \makecell{scalar quantity\\or\\vector quantity}\\
\hline
\textbf{A} & mass & scalar\\
\hline
\textbf{B} & mass & vector\\
\hline
\textbf{C} & weight & scalar\\
\hline
\textbf{D} & weight & vector\\
\hline
\end{tabular}
\end{center}

%------------------------------ 2 ------------------------------
\examq{5054/11/M/J/25}{2}{Motion or Kinematics}{1}
The speed--time graph shows part of the journey of a bicycle for a time of 10\,s.

\qfig[0.5\linewidth]{5054_11_M_J_25_fig1.png}{}

What is the distance that the bicycle travels in the 10\,s?

\vspace{2mm}
\textbf{A}\hspace{0.6em}20\,m\hspace{1.6cm}\textbf{B}\hspace{0.6em}25\,m\hspace{1.6cm}\textbf{C}\hspace{0.6em}30\,m\hspace{1.6cm}\textbf{D}\hspace{0.6em}40\,m

%------------------------------ 3 ------------------------------
\examq{5054/11/M/J/25}{3}{Forces or Dynamics}{1}
Which statement explains why a heavy coin that falls through a short distance towards the ground does not reach terminal velocity?

\begin{choices}
  \item The coin has not hit the ground.
  \item The weight of the coin equals the air resistance.
  \item The weight of the coin increases as the air resistance increases.
  \item The weight of the coin is always more than the air resistance.
\end{choices}

%------------------------------ 4 ------------------------------
\examq{5054/11/M/J/25}{4}{Deformation}{1}
The diagram shows the load--extension graph for a spring.

\qfig[0.4\linewidth]{5054_11_M_J_25_fig2.png}{}

What is the name given to point X on the graph?

\begin{choices}
  \item limit of force
  \item limit of proportionality
  \item limit of the extension
  \item limit of the spring
\end{choices}

%------------------------------ 5 ------------------------------
\examq{5054/11/M/J/25}{5}{Mass Weight and Density}{1}
Some students measure the masses and the volumes of differently sized samples of a type of wood.

Which graph shows their results?

\qfig[0.9\linewidth]{5054_11_M_J_25_fig3.png}{}

%------------------------------ 6 ------------------------------
\examq{5054/11/M/J/25}{6}{Forces or Dynamics}{1}
The diagram shows a section of a roller coaster ride that is a circle arranged vertically.

The truck is moving at constant speed at the position shown.

Which arrow shows the direction of the resultant force acting on the truck?

\qfig[0.55\linewidth]{5054_11_M_J_25_fig4.png}{}

%------------------------------ 7 ------------------------------
\examq{5054/11/M/J/25}{7}{Momentum (Dynamics)}{1}
What is a unit for momentum?

\vspace{2mm}
\textbf{A}\hspace{0.6em}kg\,m\,/\,s\hspace{1.3cm}\textbf{B}\hspace{0.6em}kg\,m\,/\,s\pow{2}\hspace{1.3cm}\textbf{C}\hspace{0.6em}N\,m\hspace{1.3cm}\textbf{D}\hspace{0.6em}N\,/\,m\pow{2}

%------------------------------ 8 ------------------------------
\examq{5054/11/M/J/25}{8}{Momentum (Dynamics)}{1}
In a safety test, a car collides with a concrete block.

The impulse exerted by the block on the car is 14\,000\,Ns.

The collision with the block lasts for 70\,ms.

What is the force exerted on the car?

\vspace{2mm}
\textbf{A}\hspace{0.6em}200\,N\hspace{1.2cm}\textbf{B}\hspace{0.6em}14\,000\,N\hspace{1.2cm}\textbf{C}\hspace{0.6em}200\,000\,N\hspace{1.2cm}\textbf{D}\hspace{0.6em}980\,000\,N

%------------------------------ 9 ------------------------------
\examq{5054/11/M/J/25}{9}{Work Energy and Power}{1}
An object pulled along horizontal ground accelerates to the right.

The diagram shows the directions of the forces acting on the object.

\qfig[0.65\linewidth]{5054_11_M_J_25_fig5.png}{}

Which force is multiplied by the distance moved to calculate the work done to increase the kinetic energy store?

\begin{choices}
  \item contact force
  \item frictional force
  \item pulling force
  \item weight
\end{choices}

%------------------------------ 10 ------------------------------
\examq{5054/11/M/J/25}{10}{Mass Weight and Density}{1}
The mass of object P is greater than the mass of object Q.

The objects contain different amounts of matter and have a different resistance to change of motion.

Which row is correct?

\begin{center}
\begin{tabular}{|c|c|c|}
\hline
 & \makecell{greater amount\\of matter} & \makecell{greater resistance\\to change of motion}\\
\hline
\textbf{A} & P & P\\
\hline
\textbf{B} & P & Q\\
\hline
\textbf{C} & Q & P\\
\hline
\textbf{D} & Q & Q\\
\hline
\end{tabular}
\end{center}

%------------------------------ 11 ------------------------------
\examq{5054/11/M/J/25}{11}{Pressure}{1}
A builder leaves two identical, heavy, stone tiles resting on soft earth. One is vertical and the other is horizontal.

After a few hours, the vertical tile has started to sink into the soft earth. The horizontal one has not started to sink.

Which row correctly compares the forces and the pressures that the tiles exert on the earth?

\begin{center}
\begin{tabular}{|c|c|c|}
\hline
 & forces & pressures\\
\hline
\textbf{A} & different & different\\
\hline
\textbf{B} & different & same\\
\hline
\textbf{C} & same & different\\
\hline
\textbf{D} & same & same\\
\hline
\end{tabular}
\end{center}

%------------------------------ 12 ------------------------------
\examq{5054/11/M/J/25}{12}{Pressure}{1}
Four beakers contain the same liquid.

At which point is the pressure the greatest?

\qfig[0.6\linewidth]{5054_11_M_J_25_fig6.png}{}

%------------------------------ 13 ------------------------------
\examq{5054/11/M/J/25}{13}{Thermal Properties \& Temperature}{1}
The strength of the forces between the particles of a solid determines how much the solid expands when heated.

Which statement about these forces is correct?

\begin{choices}
  \item The forces are strong so solids expand less than liquids and gases.
  \item The forces are strong so solids expand more than liquids and gases.
  \item The forces are weak so solids expand less than liquids and gases.
  \item The forces are weak so solids expand more than liquids and gases.
\end{choices}

%------------------------------ 14 ------------------------------
\examq{5054/11/M/J/25}{14}{Kinetic Model of Matter}{1}
Air is heated in a sealed container with constant volume.

Why does the air pressure increase when the temperature increases?

\begin{choices}
  \item The air molecules expand.
  \item The air molecules bounce off each other more frequently.
  \item The air molecules bounce off the walls more frequently.
  \item The number of air molecules increases.
\end{choices}

%------------------------------ 15 ------------------------------
\examq{5054/11/M/J/25}{15}{Thermal Properties \& Temperature}{1}
A student determines the specific heat capacity of a metal by carrying out an experiment.

The list shows the apparatus that the student has available.

\begin{center}
\begin{tabular}{@{}r@{\hspace{0.8em}}l@{}}
1 & electronic balance\\
2 & ruler\\
3 & thermometer\\
4 & heater operating at a power of 50\,W\\
5 & stopwatch\\
6 & ammeter\\
7 & voltmeter\\
\end{tabular}
\end{center}

Which apparatus does the student use?

\begin{choices}
  \item 1, 3, 4 and 5
  \item 1, 3, 6 and 7
  \item 2, 3, 4 and 6
  \item 3, 4, 5, 6 and 7
\end{choices}

%------------------------------ 16 ------------------------------
\examq{5054/11/M/J/25}{16}{Thermal Properties \& Temperature}{1}
When a liquid evaporates, molecules escape from its surface.

Which molecules escape, and what happens to the average kinetic energy of the molecules remaining in the liquid?

\begin{choices}
  \item The less energetic molecules escape and the average kinetic energy decreases.
  \item The less energetic molecules escape and the average kinetic energy increases.
  \item The more energetic molecules escape and the average kinetic energy decreases.
  \item The more energetic molecules escape and the average kinetic energy increases.
\end{choices}

%------------------------------ 17 ------------------------------
\examq{5054/11/M/J/25}{17}{Thermal Properties \& Temperature}{1}
A sample of water in a beaker is at a temperature $T$.

The surface area of the water is $A$ and an electric fan blows air across the top of the beaker at a speed $v$.

For which values of $T$, $A$ and $v$ does the water evaporate the most rapidly?

\begin{center}
\begin{tabular}{|c|c|c|c|}
\hline
 & \makecell{$T$\\\degC} & \makecell{$A$\\cm\pow{2}} & \makecell{$v$\\m\,/\,s}\\
\hline
\textbf{A} & 18 & 45 & 0.1\\
\hline
\textbf{B} & 24 & 50 & 1.9\\
\hline
\textbf{C} & 38 & 70 & 5.0\\
\hline
\textbf{D} & 38 & 45 & 4.7\\
\hline
\end{tabular}
\end{center}

%------------------------------ 18 ------------------------------
\examq{5054/11/M/J/25}{18}{Transfer of Thermal Energy}{1}
Thermal energy passes quickly from hot regions of a metal to cold regions of the metal.

Why are metals good conductors of thermal energy?

\begin{choices}
  \item Electrons in the metal vibrate about fixed points and quickly pass energy from one electron to the next.
  \item Electrons move quickly from ions in hot regions to ions in cold regions.
  \item Protons in the metal are close together and move quickly from hot to cold regions.
  \item Small ions of the metal move quickly from hot to cold regions through spaces in the metal lattice.
\end{choices}

%------------------------------ 19 ------------------------------
\examq{5054/11/M/J/25}{19}{Light (Reflection/Refraction/Lenses)}{1}
The diagram shows light striking a plane mirror and reflecting.

Which angle is the angle of incidence?

\qfig[0.5\linewidth]{5054_11_M_J_25_fig7.png}{}

%------------------------------ 20 ------------------------------
\examq{5054/11/M/J/25}{20}{Light (Reflection/Refraction/Lenses)}{1}
A shoe shop puts a mirror on the wall so that customers can look at their shoes.

The length of the mirror is 50\,cm. A customer has eyes 150\,cm above ground level.

\qfig[0.5\linewidth]{5054_11_M_J_25_fig8.png}{}

The bottom of the mirror is at height $h$ above the ground.

What is the smallest value of $h$ that allows the customer to see an image of his shoes in the mirror?

\vspace{2mm}
\textbf{A}\hspace{0.6em}0\hspace{1.6cm}\textbf{B}\hspace{0.6em}25\,cm\hspace{1.6cm}\textbf{C}\hspace{0.6em}50\,cm\hspace{1.6cm}\textbf{D}\hspace{0.6em}75\,cm

%------------------------------ 21 ------------------------------
\examq{5054/11/M/J/25}{21}{Transfer of Thermal Energy}{1}
Which description of a dull black surface is correct?

\begin{choices}
  \item good emitter, good absorber and good reflector of radiation
  \item good emitter, poor absorber and poor reflector of radiation
  \item good emitter, good absorber and poor reflector of radiation
  \item poor emitter, poor absorber and poor reflector of radiation
\end{choices}

%------------------------------ 22 ------------------------------
\examq{5054/11/M/J/25}{22}{Light (Reflection/Refraction/Lenses)}{1}
The rays of light from a ray-box pass through three lenses placed at positions 1, 2 and 3.

\qfig[0.65\linewidth]{5054_11_M_J_25_fig9.png}{}

Which type of lens is used at each position?

\begin{center}
\begin{tabular}{|c|c|c|c|}
\hline
 & position 1 & position 2 & position 3\\
\hline
\textbf{A} & converging & converging & converging\\
\hline
\textbf{B} & converging & converging & diverging\\
\hline
\textbf{C} & diverging & converging & diverging\\
\hline
\textbf{D} & diverging & diverging & converging\\
\hline
\end{tabular}
\end{center}

%------------------------------ 23 ------------------------------
\examq{5054/11/M/J/25}{23}{Light (Reflection/Refraction/Lenses)}{1}
A swimming pool is lit by an underwater light.

A ray of light is incident on the surface of the water.

What is the path for the ray of light?

\qfig[0.55\linewidth]{5054_11_M_J_25_fig10.png}{}

%------------------------------ 24 ------------------------------
\examq{5054/11/M/J/25}{24}{Light (Reflection/Refraction/Lenses)}{1}
When a converging lens is used as a magnifying glass, what is the nature of the image?

\begin{choices}
  \item real and inverted
  \item real and upright
  \item virtual and inverted
  \item virtual and upright
\end{choices}

%------------------------------ 25 ------------------------------
\examq{5054/11/M/J/25}{25}{Sound}{1}
A sound wave consists of compressions and rarefactions.

Which type of wave is sound and how does the pressure in a rarefaction compare with the pressure in a compression?

\begin{center}
\begin{tabular}{|c|c|c|}
\hline
 & type of wave & \makecell{pressure in\\a rarefaction}\\
\hline
\textbf{A} & longitudinal & smaller\\
\hline
\textbf{B} & longitudinal & greater\\
\hline
\textbf{C} & transverse & smaller\\
\hline
\textbf{D} & transverse & greater\\
\hline
\end{tabular}
\end{center}

%------------------------------ 26 ------------------------------
\examq{5054/11/M/J/25}{26}{Sound}{1}
What happens to a sound wave when the note heard gets louder?

\begin{choices}
  \item Its amplitude increases.
  \item Its frequency increases.
  \item Its speed increases.
  \item Its wavelength increases.
\end{choices}

%------------------------------ 27 ------------------------------
\examq{5054/11/M/J/25}{27}{Sound}{1}
A pulse of ultrasound from a sensor is reflected back to the sensor by a crack in a piece of metal and by the bottom surface of the metal.

\qfig[0.55\linewidth]{5054_11_M_J_25_fig11.png}{}

The reflection from the crack is received back at the sensor $1.0\times10^{-5}$\,s after the pulse is sent out.

The reflection from the bottom surface is received back at the sensor $1.5\times10^{-5}$\,s after the pulse is sent out.

The speed of ultrasound in the metal is 5200\,m\,/\,s.

What is the distance $x$ between the crack and the bottom surface of the metal?

\vspace{2mm}
\textbf{A}\hspace{0.6em}0.013\,m\hspace{1.3cm}\textbf{B}\hspace{0.6em}0.026\,m\hspace{1.3cm}\textbf{C}\hspace{0.6em}0.052\,m\hspace{1.3cm}\textbf{D}\hspace{0.6em}0.078\,m

%------------------------------ 28 ------------------------------
\examq{5054/11/M/J/25}{28}{Electrical Quantities}{1}
A 12\,V battery is connected to a resistor of resistance 100\,\ohms.

How much charge passes through the resistor in 30 minutes?

\vspace{2mm}
\textbf{A}\hspace{0.6em}3.6\,C\hspace{1.4cm}\textbf{B}\hspace{0.6em}220\,C\hspace{1.4cm}\textbf{C}\hspace{0.6em}250\,C\hspace{1.4cm}\textbf{D}\hspace{0.6em}15\,000\,C

%------------------------------ 29 ------------------------------
\examq{5054/11/M/J/25}{29}{Electrical Quantities}{1}
A circuit consists of a resistor connected between the positive terminal and the negative terminal of a power supply.

Which row describes the direction of conventional current and the direction of flow of free electrons in the resistor?

\begin{center}
\begin{tabular}{|c|c|c|}
\hline
 & \makecell{direction of\\conventional current} & \makecell{direction of flow\\of free electrons}\\
\hline
\textbf{A} & from negative to positive & from negative to positive\\
\hline
\textbf{B} & from negative to positive & from positive to negative\\
\hline
\textbf{C} & from positive to negative & from negative to positive\\
\hline
\textbf{D} & from positive to negative & from positive to negative\\
\hline
\end{tabular}
\end{center}

%------------------------------ 30 ------------------------------
\examq{5054/11/M/J/25}{30}{Electric Circuits}{1}
The graph shows how the resistance of a thermistor varies with temperature.

\qfig[0.7\linewidth]{5054_11_M_J_25_fig12.png}{}

The thermistor is at a temperature of 10\,\degC{} and is connected in series with a fixed resistor of resistance 20\,\ohms.

What is the combined resistance of both components?

\vspace{2mm}
\textbf{A}\hspace{0.6em}16\,\ohms\hspace{1.5cm}\textbf{B}\hspace{0.6em}20\,\ohms\hspace{1.5cm}\textbf{C}\hspace{0.6em}80\,\ohms\hspace{1.5cm}\textbf{D}\hspace{0.6em}100\,\ohms

%------------------------------ 31 ------------------------------
\examq{5054/11/M/J/25}{31}{Electric Circuits}{1}
Each of the resistors in the circuit shown has a resistance of 6.0\,\ohms.

\qfig[0.45\linewidth]{5054_11_M_J_25_fig13.png}{}

What is the total resistance of the circuit?

\vspace{2mm}
\textbf{A}\hspace{0.6em}2.0\,\ohms\hspace{1.5cm}\textbf{B}\hspace{0.6em}4.0\,\ohms\hspace{1.5cm}\textbf{C}\hspace{0.6em}9.0\,\ohms\hspace{1.5cm}\textbf{D}\hspace{0.6em}18\,\ohms

%------------------------------ 32 ------------------------------
\examq{5054/11/M/J/25}{32}{Electro Magnetism / Electromagnetic Induction}{1}
An alternating current (a.c.) generator is connected to an oscilloscope.

The diagram shows the trace produced on the oscilloscope screen.

\qfig[0.5\linewidth]{5054_11_M_J_25_fig14.png}{}

What are possible angles between the plane of the coil and the magnetic field direction at the times represented by points 1, 2 and 3?

\begin{center}
\begin{tabular}{|c|c|c|c|}
\hline
 & 1 & 2 & 3\\
\hline
\textbf{A} & 0\dg & 90\dg & 180\dg\\
\hline
\textbf{B} & 0\dg & 90\dg & 270\dg\\
\hline
\textbf{C} & 90\dg & 180\dg & 0\dg\\
\hline
\textbf{D} & 90\dg & 270\dg & 0\dg\\
\hline
\end{tabular}
\end{center}

%------------------------------ 33 ------------------------------
\examq{5054/11/M/J/25}{33}{Electro Magnetism / Electromagnetic Induction}{1}
A beam of beta particles travelling vertically downwards enters the magnetic field between two magnetic poles.

\qfig[0.45\linewidth]{5054_11_M_J_25_fig15.png}{}

The beta particles experience a force due to the magnetic field.

What is the direction of the force on a beta particle as it enters the magnetic field?

\begin{choices}
  \item towards the top of the page
  \item towards the bottom of the page
  \item into the page
  \item out of the page
\end{choices}

%------------------------------ 34 ------------------------------
\examq{5054/11/M/J/25}{34}{Electro Magnetism / Electromagnetic Induction}{1}
A step-up transformer is connected between a power station and a long-distance electricity transmission cable.

What is the purpose of the step-up transformer, and does it operate using alternating current (a.c.) or using direct current (d.c.)?

\begin{center}
\begin{tabular}{|c|c|c|}
\hline
 & purpose of a step-up transformer & current used\\
\hline
\textbf{A} & to decrease the current in the cable & d.c.\\
\hline
\textbf{B} & to decrease the current in the cable & a.c.\\
\hline
\textbf{C} & to increase the current in the cable & d.c.\\
\hline
\textbf{D} & to increase the current in the cable & a.c.\\
\hline
\end{tabular}
\end{center}

%------------------------------ 35 ------------------------------
\examq{5054/11/M/J/25}{35}{The Nuclear Atom}{1}
Each diagram shows the nucleus of an atom.

\qfig[0.75\linewidth]{5054_11_M_J_25_fig16.png}{}

Which diagrams show isotopes of the same element?

\vspace{2mm}
\textbf{A}\hspace{0.6em}1 and 2\hspace{1.4cm}\textbf{B}\hspace{0.6em}1 and 4\hspace{1.4cm}\textbf{C}\hspace{0.6em}2 and 3\hspace{1.4cm}\textbf{D}\hspace{0.6em}3 and 4

%------------------------------ 36 ------------------------------
\examq{5054/11/M/J/25}{36}{Radioactivity}{1}
A radioactive nucleus X decays to another radioactive nucleus. After a number of decays, a stable nucleus Q is produced.

The proton number (atomic number) of X is equal to that of Q. The nucleon number (mass number) of X is four more than that of Q.

Which radiations are emitted to produce Q?

\begin{choices}
  \item two alpha-particles
  \item one alpha-particle and one beta-particle
  \item two alpha-particles and one beta-particle
  \item one alpha-particle and two beta-particles
\end{choices}

%------------------------------ 37 ------------------------------
\examq{5054/11/M/J/25}{37}{Radioactivity}{1}
The graph shows the results of an experiment to measure the count rate of a radioactive isotope.

\qfig[0.85\linewidth]{5054_11_M_J_25_fig17.png}{}

What is the half-life of the isotope?

\begin{choices}
  \item 60 minutes
  \item 80 minutes
  \item 100 minutes
  \item 160 minutes
\end{choices}

%------------------------------ 38 ------------------------------
\examq{5054/11/M/J/25}{38}{The Nuclear Atom}{1}
During fission, a nucleus of uranium-235 (U-235) absorbs one neutron to produce two daughter nuclei and three neutrons.

Krypton-92 (Kr-92) is one of the daughter nuclei and it contains 92 nucleons.

How many nucleons does the other daughter nucleus contain?

\vspace{2mm}
\textbf{A}\hspace{0.6em}92\hspace{1.7cm}\textbf{B}\hspace{0.6em}141\hspace{1.7cm}\textbf{C}\hspace{0.6em}144\hspace{1.7cm}\textbf{D}\hspace{0.6em}327

%------------------------------ 39 ------------------------------
\examq{5054/11/M/J/25}{39}{Earth \& Solar System}{1}
The orbital speed of a planet around the Sun is $v$ and its orbital period is $T$.

Jupiter is further from the Sun than the Earth.

Which planet has the greater value for $v$ and which planet has the greater value for $T$?

\begin{center}
\begin{tabular}{|c|c|c|}
\hline
 & \makecell{greater orbital\\speed $v$} & \makecell{greater orbital\\period $T$}\\
\hline
\textbf{A} & Earth & Earth\\
\hline
\textbf{B} & Earth & Jupiter\\
\hline
\textbf{C} & Jupiter & Jupiter\\
\hline
\textbf{D} & Jupiter & Earth\\
\hline
\end{tabular}
\end{center}

%------------------------------ 40 ------------------------------
\examq{5054/11/M/J/25}{40}{Earth \& Solar System}{1}
What is the Earth's position in the order of planets from the Sun?

\begin{choices}
  \item closest planet to the Sun
  \item second closest planet to the Sun
  \item third closest planet to the Sun
  \item fourth closest planet to the Sun
\end{choices}

\end{document}